\chapter{La campaña experimental}
\label{cap:experimental}

\begin{enpocas}
Qué preguntas se quieren responder, cómo se diseñan los vuelos para responderlas sin engañarse, la
prueba de Wilcoxon con la corrección de Holm, y lo que dijo la campaña.
\end{enpocas}

\section{Las preguntas}

\begin{description}
  \item[P1.] ¿La re-observación de candidatas dudosas encuentra más personas que el barrido fijo, y
        en qué nivel de oclusión ayuda más?
  \item[P2.] ¿Cuánto cuesta? Tiempo de vuelo, personas por minuto y desvío de la ruta.
  \item[P3.] ¿Qué tan sensible es la misión a un detector más lento o que pierde cuadros
        (\cod{latencia_extra_ms}, \cod{saltar_cuadros})?
  \item[P4.] ¿Cambia algo volar con la red neuronal en vez de la cascada de PX4?
\end{description}

\section{El diseño}

Un \textbf{diseño factorial completo}: todas las combinaciones de oclusión (3), iluminación (2), viento
(2) y estrategia (2), es decir 24 vuelos por semilla con el controlador en cascada. Para P4 se agregan,
en oclusión media y sol alto, los 4 vuelos de la red (con y sin viento, las dos estrategias): 28 vuelos
por semilla.

\begin{lstlisting}[style=bash]
python3 src/mision/scripts/correr_experimentos.py --semillas 30 --salida ~/resultados     # 840 vuelos
python3 src/mision/scripts/correr_experimentos.py --semillas 2 --solo oclusion=media      # una parte
\end{lstlisting}

Cada vuelo es un \cod{ros2 launch} sin interfaz en su propia partición de Gazebo y su dominio de ROS~2,
dura unos 5~minutos y deja un JSON. El programa salta los vuelos que ya tienen JSON, así que se puede
interrumpir y continuar, o ampliar de 5 a 30 semillas.

\begin{consejo}
\textbf{Apareado.} La re-observación y el barrido fijo vuelan \emph{el mismo bosque} (misma semilla,
iluminación y viento). La diferencia entre los dos vuelos de un par se debe a la estrategia y no a que a
uno le tocó un bosque más difícil. Así se necesitan muchos menos vuelos para ver un efecto.
\end{consejo}

\section{Wilcoxon y Holm}

Para cada par se calcula la diferencia $d = \text{reobs} - \text{fija}$ (personas encontradas, por
ejemplo). La \textbf{prueba de rangos con signo de Wilcoxon} pregunta si esas diferencias tienden a ser
positivas, sin suponer que siguen una campana de Gauss: ordena las $|d|$, les pone rango y compara la
suma de los rangos positivos con la de los negativos. Se reporta la mediana de $d$ con un intervalo de
confianza del 95~\% obtenido por remuestreo (\emph{bootstrap}).

Se hacen 12 pruebas (3 oclusiones $\times$ 4 variables). Con tantas, alguna saldría ``significativa'' por
pura suerte; la \textbf{corrección de Holm} lo evita: ordena los valores $p$ de menor a mayor y multiplica
el $k$-ésimo por $(m-k+1)$, sin dejar que bajen.

\begin{lstlisting}[style=bash]
~/venv-mision/bin/python src/mision/scripts/analisis_wilcoxon.py --resultados ~/resultados
~/venv-mision/bin/python src/mision/scripts/analisis_wilcoxon.py --demo        # datos sintéticos
\end{lstlisting}

\begin{aviso}
Una sola semilla no compara estrategias: en el primer bosque la re-observación puede ganar o perder por
casualidad. Las conclusiones salen de los pares, no de un vuelo bonito.
\end{aviso}

\section{Resultados}

\IfFileExists{capitulos/resultados_campana.tex}{%
Las tablas siguientes las genera \cod{herramientas/tabla_campana.py} a partir de los JSON de
\cod{~/resultados}; se regeneran al agregar semillas.

% generado por herramientas/tabla_campana.py con 2 corridas
Corridas analizadas: 2 (semillas 1--1).
\begin{table}[H]\centering\small
\begin{tabular}{lllrrrrrrr}
\toprule
\textbf{Control} & \textbf{Oclusión} & \textbf{Estrategia} & $n$ & \textbf{Encontradas} & \textbf{FA} & \textbf{FA$_{10}$} & \textbf{Avisos} & \textbf{Duración (s)} & \textbf{RMSE (m)} \\ \midrule
cascada & baja & fija & 1 & 1/1 & 14.0 & 2.0 & 2.0 & 237 & 0.069 \\
cascada & baja & reobs & 1 & 1/1 & 22.0 & 6.0 & 5.0 & 278 & 0.279 \\
\bottomrule
\end{tabular}
\caption{Promedios por condición. FA: falsas alarmas con una detección; FA$_{10}$: con al menos 10; avisos GPS por vuelo.}
\label{tab:campana-medias}
\end{table}
\begin{table}[H]\centering\small
\begin{tabular}{llrrcrr}
\toprule
\textbf{Oclusión} & \textbf{Variable} & $n$ & \textbf{Mediana $d$} & \textbf{IC 95~\%} & $p$ & $p_{Holm}$ \\ \midrule
baja & personas & 1 & 0.000 & [0.000, 0.000] & 1.000 & 1.000 \\
baja & personas/min & 1 & -0.037 & [-0.037, -0.037] & 1.000 & 1.000 \\
baja & RMSE (m) & 1 & 0.210 & [0.210, 0.210] & 1.000 & 1.000 \\
baja & falsas alarmas & 1 & 8.000 & [8.000, 8.000] & 1.000 & 1.000 \\
media & personas & 0 & -- & [--, --] & -- & -- \\
media & personas/min & 0 & -- & [--, --] & -- & -- \\
media & RMSE (m) & 0 & -- & [--, --] & -- & -- \\
media & falsas alarmas & 0 & -- & [--, --] & -- & -- \\
alta & personas & 0 & -- & [--, --] & -- & -- \\
alta & personas/min & 0 & -- & [--, --] & -- & -- \\
alta & RMSE (m) & 0 & -- & [--, --] & -- & -- \\
alta & falsas alarmas & 0 & -- & [--, --] & -- & -- \\
\bottomrule
\end{tabular}
\caption{Re-observación menos barrido fijo, apareado por semilla, iluminación y viento (controlador en cascada).}
\label{tab:campana-wilcoxon}
\end{table}

}{%
\textit{Las tablas se generan con \cod{herramientas/tabla_campana.py} cuando termina la campaña.}
}

\begin{leccion}
\textbf{Las falsas alarmas dependen del umbral.} Con una sola detección cualquier sombra que el detector
confunda una vez cuenta como candidata (10 a 20 por vuelo); exigiendo 10 detecciones fusionadas quedan
pocas, y casi todas son grupos de robles que el detector confunde una y otra vez. Por eso la tabla trae
las dos cifras, y por eso el aviso GPS exige 10.
\end{leccion}
